% !TeX root = ../../Book5.tex
% 纲要标题：### 7.4 张量分解
% 纲要位置：工作记录/计划纲要.md，第 4295 行。
% * Schur 函子
% * Pieri 规则及低阶 Littlewood–Richardson 例子
% * Littlewood–Richardson 诱导分解先在有理数上确定重数，再扩张到任意特征零域；一般线性群分解由同一个限制表示的重数空间及 Frobenius 互反确定
% * 与第二卷对称幂、外幂的统一，以及 \((2,1)\) 型的乘法核模型
% ---
\section{张量分解}
\label{b5:sec:0704}

Schur--Weyl 对偶给每个分拆配置了一个重数空间。改变向量空间或施加线性映射时，这些空间也会相容地改变，因而形成函子。对称函数的乘法随后给出其张量积规则；需要注意，这种乘法对应对称群一侧的诱导，不能误认为同一对称群上的内部张量积。

\subsection{Schur 函子}
\label{b5:subsec:0704:01}

\begin{definition}\label{b5:def:schur-functor}
设 \(k\) 为特征零域，\(n\) 为非负整数，\(\lambda\vdash n\)，固定一个 Specht 左 \(k[S_n]\)-模 \(S^\lambda_k\)。对有限维 \(k\) 向量空间 \(V\)，在 \(V^{\otimes n}\) 上采用置换张量因子的左 \(S_n\) 作用，令
\[
\mathbb S_\lambda(V)
=\operatorname{Hom}_{k[S_n]}(S^\lambda_k,V^{\otimes n}).
\]
对有限维 \(k\) 向量空间之间的线性映射 \(f:V\to W\)，以及 \(u\in\mathbb S_\lambda(V)\)，令
\[
\mathbb S_\lambda(f)(u)=f^{\otimes n}\circ u.
\]
这称为分拆 \(\lambda\) 对应的\term[Schur-hanzi]{Schur 函子}[Schur functor]。空分拆对应常值空间 \(k\)；正次数时规定 \(\mathbb S_\lambda(0)=0\)。
\end{definition}
\symidx{\(\mathbb S_\lambda\)：Schur 函子}

因为 \(f^{\otimes n}\) 与因子置换相容，上述复合仍为 \(S_n\) 同态。恒等映射和复合也被保持，所以确实得到函子。选定基时，其作用由 \(f\) 的矩阵条目的 \(n\) 次齐次多项式给出；通常它不是关于 \(f\) 的线性运算。例如 \(\operatorname{Sym}^2(f+g)\) 中有交叉项，一般不等于 \(\operatorname{Sym}^2(f)+\operatorname{Sym}^2(g)\)。

\ProseParagraphBreak
这个定义中的空间就是上一节的 \(L_\lambda(V)\)。因此它在非零时是不可约的一般线性群表示，其迹函数由 Schur 多项式 \(s_\lambda\) 给出：把群元素在代数闭包中的特征值代入即可求迹。分拆的行数超过 \(\dim V\) 时，这个空间为零。在不同维数下使用同一分拆，不需要重新定义构造。

\ProseParagraphBreak

也可以把这个 Hom 空间具体识别为 Young 对称化子的像。对 \(n\ge1\)，选定形状为 \(\lambda\) 的表 \(t\)，令
\[
A=k[S_n],\qquad T=V^{\otimes n},\qquad
e=\frac{f^\lambda}{n!}b_ta_t.
\]
由\cref{b5:prop:young-quasi-idempotent}，\(e^2=e\)，且选定的左模同构 \(S_k^\lambda\cong Ae\) 将 Schur 函子识别为 \(\operatorname{Hom}_A(Ae,T)\)。这时有双向映射
\[
\begin{aligned}
\operatorname{Hom}_A(Ae,T)&\longrightarrow eT,
&F&\longmapsto F(e),\\
eT&\longrightarrow\operatorname{Hom}_A(Ae,T),
&v&\longmapsto\bigl(ae\longmapsto av\bigr).
\end{aligned}
\]
第一式落在 \(eT\) 中，因为 \(eF(e)=F(e^2)=F(e)\)。第二式良定义：若 \(ae=a'e\)，则
\((a-a')v=(a-a')ev=0\)。它显然为左 \(A\)-模同态，而 \(F(ae)=aF(e)\) 与 \(ev=v\) 分别表明两个复合为恒等。对 \(g\in\operatorname{GL}_k(V)\)，与因子置换交换的 \(g^{\otimes n}\) 在两侧给出相容作用。更一般地，线性映射 \(f:V\to W\) 诱导的 \(f^{\otimes n}\) 是 \(A\)-模同态，故把 \(eV^{\otimes n}\) 送入 \(eW^{\otimes n}\)，并满足
\[
(f^{\otimes n}\circ F)(e)=f^{\otimes n}(F(e)).
\]
这说明“在 \(e\) 处求值”与 \(f\) 的作用相容；其逆识别也相容，因此这次识别随 \(V\) 自然变化。

\ProseParagraphBreak

这里的自然性是指固定 \(t\) 后与线性映射 \(V\to W\) 相容，并不是说不同的表在同一个张量幂中给出相同的子空间。例如取 \(\lambda=(2,1)\)，令 \(t=12|3\)、\(t'=13|2\)，并记它们归一化后的 Young 对称化子为
\[
e_t=\frac13(1-(13))(1+(12)),\qquad
e_{t'}=\frac13(1-(12))(1+(13)).
\]
在\cref{b5:exm:specht-21}的三维表类置换模中，\(v_1,v_2,v_3\) 是彼此独立的表类基；其 Specht 子模取基 \(u=v_3-v_1\)、\(v=v_2-v_1\)。先作行求和，再作列交错，有
\[
\begin{aligned}
(1+(12))u&=2v_3-v_1-v_2,&
(1-(13))(2v_3-v_1-v_2)&=3u,\\
(1+(13))v&=2v_2-v_1-v_3,&
(1-(12))(2v_2-v_1-v_3)&=3v.
\end{aligned}
\]
同时 \((1+(12))v=0\)、\((1+(13))u=0\)。除以 \(3\) 即得
\[
e_tu=u,\quad e_tv=0,\qquad
e_{t'}u=0,\quad e_{t'}v=v.
\]
因此在 \(V=k^2\) 的等型分量
\(S_k^{(2,1)}\otimes L_{(2,1)}(V)\) 中，两者的像分别是
\(ku\otimes L_{(2,1)}(V)\) 与 \(kv\otimes L_{(2,1)}(V)\)。它们是不同的二维子空间，却都给出同一个 Schur 函子的具体模型；整个等型分量则为四维。这里须区分一个选定的 Young 像、抽象的重数空间，以及包含所有副本的等型分量。

\ProseParagraphBreak

这也给出行数限制的直接原因。\(b_t\) 包含每列的反对称求和；若第一列的长度 \(\ell(\lambda)\) 超过 \(\dim V\)，该列因子上的反对称张量空间就是 \(\bigwedge^{\ell(\lambda)}V=0\)。特征零下反对称求和是列长阶乘乘以反对称投影，故 \(b_tT=0\)，从而 \(eT=0\)。抽象的 Hom 空间与列反对称化的消失于是描述同一件事。

\begin{definition}\label{b5:def:lr-coefficient}
设 \(a,b\) 为非负整数，\(\lambda\vdash a\)、\(\mu\vdash b\)。在有理系数对称函数代数 \(\Lambda_{\mathbb Q}\) 中，用 Schur 基展开乘积：
\[
s_\lambda s_\mu=\sum_{\nu\vdash a+b}c_{\lambda\mu}^{\nu}s_\nu.
\]
系数 \(c_{\lambda\mu}^{\nu}\) 称为\term[Littlewood-Richardson-xishu]{Littlewood--Richardson 系数}[Littlewood--Richardson coefficient]。\footnote{利特尔伍德（Dudley Ernest Littlewood，1903-1979）与理查森（Archibald Read Richardson，1881-1954）均为英国数学家，合作研究对称群表示与对称函数的乘法系数。}
\end{definition}
\symidx{\(c_{\lambda\mu}^{\nu}\)：Littlewood--Richardson 系数}

\begin{proposition}\label{b5:prop:schur-tensor-coefficients}
设 \(k\) 为特征零域，\(a,b\) 为非负整数，\(\lambda\vdash a\)、\(\mu\vdash b\)，\(c_{\lambda\mu}^{\nu}\) 为 Schur 函数乘积的 Littlewood--Richardson 系数，\(S^\eta\) 表示 \(k\) 上相应的 Specht 左模，\(\boxtimes\) 表示外张量积，\(S_a\times S_b\le S_{a+b}\) 按两个连续字母块嵌入。则所有 \(c_{\lambda\mu}^{\nu}\) 是非负整数，且
\[
\operatorname{Ind}_{S_a\times S_b}^{S_{a+b}}
(S^\lambda\boxtimes S^\mu)
\simeq\bigoplus_{\nu\vdash a+b}(S^\nu)^{\oplus c_{\lambda\mu}^{\nu}}.
\]
对任意有限维 \(k\) 向量空间 \(V\)，以 \(\mathbb S_\eta\) 记 Schur 函子，还有 \(\operatorname{GL}_k(V)\)-表示的同构
\begin{equation}\label{b5:eq:schur-tensor-decomposition}
\mathbb S_\lambda(V)\otimes\mathbb S_\mu(V)
\simeq
\bigoplus_{\substack{\nu\vdash a+b\\\ell(\nu)\le\dim V}}
\mathbb S_\nu(V)^{\oplus c_{\lambda\mu}^{\nu}}.
\end{equation}
\end{proposition}
\begin{proof}
若 \(a+b=0\)，则 \(\lambda,\mu\) 都是空分拆，唯一的乘积系数为一，两式两边均为 \(k\)。以下设 \(a+b>0\)。先在复数域上，由诱导乘积公式及 Specht--Schur 对应得到第一式；其系数是实际分解重数，所以是非负整数。令
\[
I_{\mathbb Q}=\operatorname{Ind}_{S_a\times S_b}^{S_{a+b}}
\bigl(S_{\mathbb Q}^\lambda\boxtimes S_{\mathbb Q}^\mu\bigr).
\]
由\eqref{b5:eq:symmetric-split-algebra}，\(\mathbb Q[S_{a+b}]\) 是分裂半单代数，故
\[
I_{\mathbb Q}\cong\bigoplus_{\nu\vdash a+b}
(S_{\mathbb Q}^\nu)^{\oplus m_\nu},
\qquad m_\nu\in\mathbb Z_{\geq0}.
\]
整数标准基给出 \(k\otimes_{\mathbb Q}S_{\mathbb Q}^\eta\cong S_k^\eta\)。在群代数张量模型中结合张量积，又得到诱导与扩张标量相容的同构
\[
k\otimes_{\mathbb Q}I_{\mathbb Q}
\cong\operatorname{Ind}_{S_a\times S_b}^{S_{a+b}}
\bigl(S_k^\lambda\boxtimes S_k^\mu\bigr).
\]
先取 \(k=\mathbb C\)。各个 \(S_{\mathbb Q}^\nu\) 复化后仍是互不同构的简单模，比较复数域上的分解即得 \(m_\nu=c_{\lambda\mu}^{\nu}\)。再取任意特征零域 \(k\)，上述同构给出第一式。

对第二式，记 \(H=S_a\times S_b\)、\(U=S_k^\lambda\boxtimes S_k^\mu\)、\(T=V^{\otimes(a+b)}\)。考虑同一个重数空间
\[
Q=\operatorname{Hom}_{k[H]}(U,T),
\]
一般线性群通过在 \(T\) 上的作用作用于 \(Q\)。先在整个 \(S_{a+b}\) 下分解 \(T\)，再限制到 \(H\)，得到
\[
Q\cong\bigoplus_{\nu\vdash a+b}
\operatorname{Hom}_{k[H]}\!\left(U,\operatorname{Res}_{H}^{S_{a+b}}S_k^\nu\right)
\otimes_k\mathbb S_\nu(V).
\]
由\cref{b5:thm:frobenius-reciprocity}，前面的 Hom 空间同构于
\(\operatorname{Hom}_{k[S_{a+b}]}(\operatorname{Ind}_{H}^{S_{a+b}}U,S_k^\nu)\)。第一式已经给出诱导模的分解，而每个 Specht 模的自同态只有标量，所以这个空间的维数恰为 \(c_{\lambda\mu}^{\nu}\)。

另一方面，按前 \(a\) 个、后 \(b\) 个因子重组张量幂，并分别使用 Schur--Weyl 分解，有
\[
T\cong\bigoplus_{\alpha\vdash a,\,\beta\vdash b}
(S_k^\alpha\boxtimes S_k^\beta)
\otimes_k\bigl(\mathbb S_\alpha(V)\otimes_k\mathbb S_\beta(V)\bigr).
\]
不同外张量积是 \(k[H]\) 的互不同构简单模，它们的自同态环均为 \(k\)。从这个分解取 \(U\) 的 Hom，只留下 \((\alpha,\beta)=(\lambda,\mu)\)，于是
\[
Q\cong\mathbb S_\lambda(V)\otimes_k\mathbb S_\mu(V).
\]
比较这两个对 \(Q\) 的描述就得到第二式，行数过多的 Schur 函子为零，故删去。次数零的因子采用空分拆与常值空间 \(k\)，上述论证仍成立；若 \(V=0\)，因 \(a+b>0\)，两边都是零。这也说明两条分解公式为何使用同一组系数：它们都在计算限制表示中同一个 \(U\) 的重数空间。对线性映射 \(f:V\to W\)，张量幂 \(f^{\otimes(a+b)}\) 与各因子置换交换，因而在上述重数空间上诱导 \(\theta\mapsto f^{\otimes(a+b)}\circ\theta\)。两个分解都由评价映射 \(s\otimes\theta\mapsto\theta(s)\) 得到，而先评价再取张量幂与先复合再评价相同。因此这些对应与 \(f\) 相容。最后选取只涉及有限群表示、与 \(V\) 无关的 Hom 空间基，就把固定的重数空间写成若干份 \(k\)，仍保留这种相容性。
\end{proof}

上述论证解释了分解系数的表示意义，一般 Littlewood--Richardson 表格规则尚需另作证明。下面证明最常用的单行、单列情形，再计算一个出现重数二的例子。

\subsection{Pieri 规则与 Littlewood–Richardson 低阶例子}
\label{b5:subsec:0704:02}

\begin{definition}\label{b5:def:horizontal-vertical-strip}
设分拆 \(\mu\) 的 Young 图包含分拆 \(\lambda\) 的 Young 图。差图 \(\mu/\lambda\) 若每列至多有一格，称为\term[hengtiao]{横条}[horizontal strip]；若每行至多有一格，称为\term[shutiao]{竖条}[vertical strip]。其大小是新增格数 \(|\mu|-|\lambda|\)。
\end{definition}

下面的 Pieri 规则\footnote{皮耶里（Mario Pieri，1860-1913），意大利数学家，研究几何基础，其名字也见于对称函数的乘法规则。}用横条、竖条描述 Schur 函数分别乘以单行、单列函数的结果。

\begin{theorem}[Pieri 规则]\label{b5:thm:pieri}
在有理系数对称函数代数 \(\Lambda_{\mathbb Q}\) 中，记 \(s_\lambda\) 为 Schur 函数、\(h_r,e_r\) 分别为完全、初等对称函数。对分拆 \(\lambda\) 和整数 \(r\ge0\)，有
\begin{align}
s_\lambda h_r&=\sum_{\mu/\lambda\text{ 为大小 }r\text{ 的横条}}s_\mu,
\label{b5:eq:pieri-horizontal}\\
s_\lambda e_r&=\sum_{\mu/\lambda\text{ 为大小 }r\text{ 的竖条}}s_\mu.
\label{b5:eq:pieri-vertical}
\end{align}
每项的系数均为一。
\end{theorem}
\begin{proof}
取 \(N\ge|\lambda|+r\)，令 \(\alpha=\lambda+\delta_N\)。先证横条公式。由\cref{b5:thm:jacobi-trudi}证明中按单项式轨道得到的交错多项式基，\(a_\alpha h_r\) 可唯一写成严格递减非负指数列 \(\beta\) 对应的 \(a_\beta\) 的线性组合。\(a_\beta\) 中单项式 \(x^\beta\) 的系数为一，而其他基元素中该系数为零，所以展开系数正是 \(a_\alpha h_r\) 中 \(x^\beta\) 的系数。由于 \(h_r\) 中每个次数 \(r\) 的单项式系数都是一，该系数为
\[
\sum_{\sigma\in S_N}\operatorname{sgn}(\sigma)
\prod_{i=1}^N\mathbf1_{\bigl\{\beta_i\ge\alpha_{\sigma(i)}\bigr\}}
=\det\bigl(\mathbf1_{\{\beta_i\ge\alpha_j\}}\bigr)_{i,j=1}^N,
\]
其中已限定 \(|\beta|=|\alpha|+r\)。

\(\alpha,\beta\) 都严格递减。因此矩阵第 \(i\) 行先为若干个零，随后全为一；其第一个一的位置 \(q_i\) 随 \(i\) 不减。如果有零行或重复行，行列式为零。其余情形的 \(N\) 个不同起始位置只能依次为 \(1,\ldots,N\)，即 \(q_i=i\)，矩阵上三角且对角元全为一。这个条件等价于
\[
\beta_1\ge\alpha_1,\qquad
\alpha_i\le\beta_i<\alpha_{i-1}\quad(2\le i\le N).
\]
写 \(\beta=\mu+\delta_N\)，它化为
\[
\mu_i\ge\lambda_i,\qquad
\mu_i\le\lambda_{i-1}\quad(i\ge2).
\]
恰好表示新增格子不能在同一列重复，即横条条件。除以 \(\Delta_N\)，得到第一式。

例如取 \(\lambda=(2)\)、\(r=2\)、\(N=4\)，则 \(\alpha=(5,2,1,0)\)。对候选形状 \(\mu=(2,1,1)\) 和 \(\mu=(3,1)\)，相应的 \(\beta\) 分别为 \((5,3,2,0)\) 与 \((6,3,1,0)\)，上述零一矩阵分别是
\[
\begin{pmatrix}
1&1&1&1\\0&1&1&1\\0&1&1&1\\0&0&0&1
\end{pmatrix},\qquad
\begin{pmatrix}
1&1&1&1\\0&1&1&1\\0&0&1&1\\0&0&0&1
\end{pmatrix}.
\]
前者的第二、三行相同，行列式为零；新增的两格都在第一列，确实不是横条。后者上三角且对角元为一，行列式为一；新增格位于第一行第三列和第二行第一列，确实是横条。这两个矩阵把“列中不能重复添格”转成了行列式中的相消与保留。稳定检测保证有限变量证明给出第一式在 \(\Lambda_{\mathbb Q}\) 中的恒等式。

再用\cref{b5:prop:omega-sign-conjugation}的对合 \(\omega\) 证明竖条公式。对 \(\lambda'\) 的横条公式施加 \(\omega\)，左边变成 \(s_\lambda e_r\)，右边各项变成 \(s_{\mu'}\)。转置 Young 图把每列至多一格变成每行至多一格，所以 \(\mu/\lambda'\) 是横条，当且仅当 \(\mu'/\lambda\) 是竖条。把 \(\mu'\) 重命名，便得到第二式，且每项系数仍为一。
\end{proof}

例如
\[
s_{(2,1)}h_2=s_{(4,1)}+s_{(3,2)}+s_{(3,1,1)}+s_{(2,2,1)}.
\]
\((3,2)\) 的新增格位于第一行第三列和第二行第二列，属于横条；不能把“横条”误解为全部新格须在同一行。

\begin{example}\label{b5:exm:lr-square-21}
利用 \(s_{(2,1)}=h_2h_1-h_3\)，计算 \(s_{(2,1)}^2\)。先按横条规则乘 \(h_2\)，得到上面的四项，再逐项添一格：
\begin{align*}
s_{(2,1)}h_2h_1
={}&s_{(5,1)}+2s_{(4,2)}+2s_{(4,1,1)}+s_{(3,3)}
+3s_{(3,2,1)}\\
&+s_{(3,1,1,1)}+s_{(2,2,2)}+s_{(2,2,1,1)}.
\end{align*}
另一方面，
\[
s_{(2,1)}h_3
=s_{(5,1)}+s_{(4,2)}+s_{(4,1,1)}+s_{(3,2,1)}.
\]
相减得到
\begin{align}
s_{(2,1)}^2
={}&s_{(4,2)}+s_{(4,1,1)}+s_{(3,3)}+2s_{(3,2,1)}
\notag\\
&+s_{(3,1,1,1)}+s_{(2,2,2)}+s_{(2,2,1,1)}.
\label{b5:eq:lr-square-21}
\end{align}
因此一般 Littlewood--Richardson 系数不一定是零或一。对三维空间，只保留至多三行的项；它们的维数依次为 \(27,10,10,8,1\)，其中八维项出现两次，合计 \(64\)，与 \(\dim\mathbb S_{(2,1)}(k^3)^2=8^2\) 相符。
\end{example}

\subsection{Schur 函子与对称幂、外幂}
\label{b5:subsec:0704:03}

\begin{proposition}\label{b5:prop:schur-symmetric-exterior}
设 \(k\) 为特征零域，\(V\) 为有限维 \(k\) 向量空间，\(n\ge1\) 为整数，\(\mathbb S_\lambda\) 为分拆 \(\lambda\) 对应的 Schur 函子。则自然地有
\[
\mathbb S_{(n)}(V)\simeq\operatorname{Sym}^nV,\qquad
\mathbb S_{(1^n)}(V)\simeq\bigwedge^nV.
\]
因而对任意分拆 \(\lambda\) 与非负整数 \(r\)，张量积
\(\mathbb S_\lambda(V)\otimes\operatorname{Sym}^rV\) 按横条分解，
\(\mathbb S_\lambda(V)\otimes\bigwedge^rV\) 按竖条分解；行数超过 \(\dim V\) 的项删去。
\end{proposition}
\begin{proof}
\(S^{(n)}\) 为平凡模，故 \(\mathbb S_{(n)}(V)\) 是对称张量子空间。对称幂的商映射 \(q:V^{\otimes n}\to\operatorname{Sym}^nV\)，在该子空间上是同构，其逆为
\[
v_1\cdots v_n\longmapsto
\frac1{n!}\sum_{\sigma\in S_n}
v_{\sigma(1)}\otimes\cdots\otimes v_{\sigma(n)}.
\]
右式对因子排列不变，所以良定；与商映射复合为恒等，对已对称张量再平均也为恒等。故两个复合都正确，且与线性映射相容。

\(S^{(1^n)}\) 为符号模，故相应 Hom 空间识别为反对称张量。外幂的商映射在此子空间上的逆为
\[
v_1\wedge\cdots\wedge v_n\longmapsto
\frac1{n!}\sum_{\sigma\in S_n}\operatorname{sgn}(\sigma)
v_{\sigma(1)}\otimes\cdots\otimes v_{\sigma(n)}.
\]
有两个相等向量时，右端按交换它们的置换配对相消，所以通过外幂的定义关系；两次复合同样为恒等。最后把这两种识别代入\eqref{b5:eq:schur-tensor-decomposition}，并使用\cref{b5:thm:pieri}，即得所述分解。
\end{proof}

例如
\[
\operatorname{Sym}^2V\otimes V
\simeq\operatorname{Sym}^3V\oplus\mathbb S_{(2,1)}(V),
\qquad
\bigwedge^2V\otimes V
\simeq\mathbb S_{(2,1)}(V)\oplus\bigwedge^3V.
\]
因此同一个 \(\mathbb S_{(2,1)}\) 同时出现在“先对称两个因子”和“先反对称两个因子”这两种构造中。对称、反对称只是一般 Young 对称性的两端情形。

\begin{proposition}\label{b5:prop:schur-21-kernels}
设 \(k\) 为特征零域，\(V\) 为有限维 \(k\) 向量空间。对称代数的乘法与外代数的乘法分别给出满射
\[
\begin{aligned}
m:\operatorname{Sym}^2V\otimes V&\longrightarrow\operatorname{Sym}^3V,
&(uv)\otimes w&\longmapsto uvw,\\
a:\bigwedge^2V\otimes V&\longrightarrow\bigwedge^3V,
&(u\wedge v)\otimes w&\longmapsto u\wedge v\wedge w.
\end{aligned}
\]
它们与 \(V\) 的线性映射相容，且其核都与 \(\mathbb S_{(2,1)}(V)\) 同构。这些同构可取为随 \(V\) 自然变化的同构。
\end{proposition}
\begin{proof}
上述映射由相应代数的乘法定义，因此良定义；三次乘积生成各自的三次齐次部分，所以满射。它们有与线性映射相容的右逆
\[
\begin{aligned}
s(uvw)&=\tfrac13\bigl((uv)\otimes w+(uw)\otimes v+(vw)\otimes u\bigr),\\
j(u\wedge v\wedge w)&=\tfrac13\bigl((u\wedge v)\otimes w
+(v\wedge w)\otimes u+(w\wedge u)\otimes v\bigr).
\end{aligned}
\]
第一式对三个向量对称，第二式对三个向量交错，因此分别通过对称幂和外幂的定义关系。再乘回去，三项各给出原乘积的三分之一，故 \(ms=1\)、\(aj=1\)。

在前面的 Pieri 分解中，\(\operatorname{Sym}^3V\) 与 \(\bigwedge^3V\) 各出现一次，其余项只有 \(\mathbb S_{(2,1)}(V)\)。非零的不同形状项是互不同构的不可约一般线性群表示，所以 \(m\)、\(a\) 在 \((2,1)\) 项上均为零，在其目标类型上为同构。因此两个核正好是 \((2,1)\) 项；为零的项不影响这个结论。\Cref{b5:prop:schur-tensor-coefficients}的证明将分解同构取为随 \(V\) 自然的同构，故这两次核的识别也可取为自然的。
\end{proof}

例如 \(V=k^3\) 时，两个核的维数分别为
\[
6\cdot3-10=8,\qquad 3\cdot3-1=8,
\]
都与 \(\dim\mathbb S_{(2,1)}(k^3)=8\) 一致。这样，\((2,1)\) 型不只由抽象 Hom 定义，也能在熟悉的对称幂、外幂及其乘法中具体找到。

\ProseParagraphBreak
这些分解由有限群与张量计算得到。研究一般线性 Lie 代数时，分拆还会作为最高权出现，从而给出另一种解释。Schur 多项式的迹与维数计算可对照\cite[Sections 5.21--5.22, pp. 128--130]{EtingofEtAl2011RepresentationTheory}。
